% Figure: beam-column interaction diagram. Axial ratio on the vertical
% axis, moment ratio on the horizontal. Shows the AISC bilinear envelope
% (steeper slope for P/Pc >= 0.2, shallower below) and a sample demand
% point inside the safe region.
\begin{figure}[htbp]
\centering
\begin{tikzpicture}
\begin{axis}[
  width=8.4cm, height=7.2cm,
  xlabel={moment ratio $M_r/M_c$},
  ylabel={axial ratio $P_r/P_c$},
  xmin=0, xmax=1.15, ymin=0, ymax=1.15,
  xtick={0,0.2,0.4,0.6,0.8,1.0},
  ytick={0,0.2,0.4,0.6,0.8,1.0},
  axis lines=left,
  tick label style={font=\scriptsize},
  label style={font=\scriptsize},
  legend style={font=\scriptsize, at={(0.97,0.97)}, anchor=north east},
]
% AISC H1-1 bilinear envelope:
% for P/Pc >= 0.2: P/Pc + (8/9)(M/Mc) = 1
% for P/Pc <  0.2: P/Pc/2 + (M/Mc) = 1
% upper segment: from (M=0,P=1) to the kink at P=0.2
% at P=0.2: M/Mc = (1-0.2)*9/8 = 0.9
\addplot[engblue, very thick] coordinates {(0,1.0) (0.9,0.2)};
% lower segment: from kink (0.9,0.2) to (M=1, P=0)
% check: P/Pc/2 + M/Mc = 1 -> at P=0.2: 0.1 + M =1 -> M=0.9 OK; at P=0 -> M=1
\addplot[engblue, very thick] coordinates {(0.9,0.2) (1.0,0.0)};
\addlegendentry{AISC envelope}
% safe-region shading
\addplot[engblue!10, draw=none] coordinates {(0,0) (0,1.0) (0.9,0.2) (1.0,0.0)} \closedcycle;
% sample demand point inside
\addplot[only marks, mark=*, engred, mark size=2pt] coordinates {(0.45,0.35)};
\node[engred, anchor=west, font=\scriptsize] at (axis cs:0.47,0.37) {demand};
% kink marker
\addplot[only marks, mark=*, engblack, mark size=1.2pt] coordinates {(0.9,0.2)};
\node[engblack, anchor=south west, font=\scriptsize] at (axis cs:0.82,0.22) {kink at $0.2$};
\end{axis}
\end{tikzpicture}
\caption[Beam-column interaction diagram]{The interaction envelope for
a member carrying combined axial compression and bending. The vertical
axis is the required axial force as a fraction of the axial capacity,
the horizontal axis the required moment as a fraction of the moment
capacity. The bilinear envelope steepens above an axial ratio of
$0.2$, where the axial term dominates, and flattens below it, where
bending dominates. Any demand point inside the envelope passes the
combined check; a point outside fails even though each load taken alone
might be acceptable. The kink reflects the design judgement that the
two failure modes interact differently depending on which load is the
larger.}
\label{fig:vol03ch08:beam-column-interaction}
\end{figure}
